\section{模型评价与展望}
\subsection{模型优点}

\begin{enumerate}
\item \textbf{优点一：以官方全局完成时间闭合建模与求解。} 三问都把切图、同核顺序、跨核依赖和存储行为落实为可执行方案，再由题定评估器计算 Makespan。候选的粗负载、连通块大小或命中字节率只是构造和诊断量，不被误写成实际完成时间。问题一的五核平均加速比为$3.5093$，问题二为$3.5456$，问题三为$3.6104$；这些数值均来自100个官方算例的完整配置，而非代理目标的预测。

\item \textbf{优点二：使用成对反事实分离场景与缓存机制。} 问题二把问题一的同一方案重新放入场景B，先测量Task合并、私有存储驻留与COPY规则改变的净影响，再与场景B重新择优的附加收益区分。问题三在相同算例、核数、切图和同核顺序上比较无L2与题定L2评估结果，使场景差异与候选更替不混入Cache机制收益。逐例比例而非只有总均值保留了有益、无益和不利的真实差异。

\item \textbf{优点三：保留可追溯的方案、响应和独立复核。} 对100例与五种核数逐一保存方案、官方原始响应和结构化摘要；问题一至问题三所选配置各有500份独立检查记录。校验覆盖操作唯一性、子图依赖、同核顺序及关键时间线，减少将评估器成功退出误当作科学验证的风险。候选失败也被保留在模型比较记录中，不用缺失分数冒充劣势。

\item \textbf{优点四：把结构异质性纳入解释。} 输入图的弱连通块与宽度差异决定连通块、深度分段和分批候选的适用范围；正式结果同时报告分位数、逐例非单调性与候选胜出频数。这样，平均加速比之外的低收益尾部及新增DDR搬运都能与调度机制对应，而不被单一汇总指标掩盖。
\end{enumerate}

\subsection{模型缺点}

\begin{enumerate}
\item \textbf{缺点一：有限候选集不保证全局最优。} 方案只在预先确定的连通块、八段深度划分和宽图分批三类候选中择优，没有穷举合法切图与核内顺序。特别是较大的弱连通块若需要非均匀切分，固定深度段可能错过更好的并行安排。因此文中加速比是可复核的可行结果，不能宣称达到理论最优。

\item \textbf{缺点二：粗负载未精确预报资源争用。} 连通块权重主要汇总操作计算周期，对共享DDR带宽、四条Pipe、固定Task等待及L1/UB容量的联合作用没有建立闭式代理。它足以产生合法候选，却可能把局部均衡方案送入全局拥塞；正式结果中随核数增加仍出现逐例速度下降，说明代理与Makespan之间存在结构依赖。

\item \textbf{缺点三：问题三的参数外推范围有限。} 正式成绩只对应题定L2容量和带宽。另选六个结构不同的算例做54格容量—带宽网格，是解释命中和瓶颈的非官方机制实验，并非官方100例的重新计分。网格中候选身份保持不变，能帮助隔离参数影响，却不能代表所有算例在其他硬件配置下的最优策略。

\item \textbf{缺点四：模型依赖题给评估语义。} FIFO不刷新次序、COPY\_IN按逻辑张量编号查询和两条读带宽池等规则都按官方实现核验。真实处理器若增加写回、异步预取、不同缓存一致性或变动的Task启动成本，原候选排序可能改变，须重新标定而不能直接套用本文速度比。
\end{enumerate}

\subsection{模型的推广与改进}

这套“依赖图切分—可执行方案—事件评估—成对反事实”的程序可用于具有明确操作依赖、可配置并行执行资源和可观察搬运成本的其他推理或数据流处理器。推广的前提是准确重建目标设备的合法性规则、存储层级、Pipe并发与固定开销，并提供与之匹配的评估或实机计时接口。若只有粗略峰值算力而缺少资源争用数据，本文的方案生成思想仍可借用，但不能迁移本文的数值加速结论。

针对有限候选集，后续可在保持覆盖与无环约束的前提下，对大连通块增加按关键路径、边界张量字节和容量余量联合引导的自适应切分；对已有方案实施局部合并、迁核和同核重排。每次改动应在固定案例划分和相同官方评估预算下成对比较，同时报告改进幅度与额外计算成本。针对缓存参数，宜扩大结构分层抽样，并在不同容量与带宽上重新优化候选，与同方案机制反事实分别报告；只有当独立复核和完整输入范围保持成立，才可将其作为新的正式策略。
